% Reorganized shared appendix file; load once in the documented order.
\section{A basic toolbox for siphon calculus}
\label{ss:siphonToolbox}

Throughout this toolbox, $\dot x=f(x)$ is a positive $C^1$ system on a
neighbourhood of the orthant; local Lipschitz regularity suffices for
invariance and mass balance. We use the face, block and spectral notation
of Section~\ref{ss:siphons}. Reaction-network statements retain the
exact-zero rate condition of Proposition~\ref{p:zeroset}.

\paragraph{Conservation classes and inaccessible extinction faces.}
Let $S$ be the stoichiometric matrix. A nonnegative conservation vector
means $q\geq0$, $q\ne0$, $q^\top S=0$; its associated conserved quantity
is $q^\top x$. Conservation laws with mixed signs are not substitutes
for this condition.

\beP[Conservation excludes a siphon face]\label{p:toolConservation}
Suppose $q\geq0$, $q\ne0$, $q^\top S=0$, and
$\operatorname{supp}q\subseteq\Sigma$. On a trajectory with
$q^\top x(0)=c>0$,
\[
 q_\Sigma^\top x_\Sigma(t)=c,
 \qquad \|x_\Sigma(t)\|_1\geq c/\|q\|_\infty.
\]
Consequently that trajectory cannot approach $\mathcal F_\Sigma$.
For a reaction network satisfying the exact-zero condition, if every
nonempty siphon contains such a conservation support, then every bounded
trajectory starting strictly positive satisfies
$\liminf_{t\to\infty}x_i(t)>0$ for every $i$.
\eeP
\begin{proof}
The first assertion follows from $q^\top\dot x=0$ and
$q_\Sigma^\top x_\Sigma\leq\|q\|_\infty\|x_\Sigma\|_1$.
For the last assertion, a bounded trajectory has a compact omega-limit
set. By Proposition~\ref{p:zeroset}, the nonempty exact zero set of a
boundary limit point would be a siphon. Its conservation support has
positive conserved mass along the trajectory but zero mass at that point,
a contradiction. Compactness gives the stated coordinate lower bounds.
\end{proof}
This is the familiar noncritical-siphon obstruction to extinction
\cite{AdLS,ShiuStu}. It yields trajectory-dependent lower bounds;
it does not assert one uniform bound across conservation classes whose
conserved masses may approach zero. A conserved quantity can also
produce zero eigenvalues in the ambient Jacobian: stability must then
be tested within the stoichiometric compatibility class.
\iffalse
\paragraph{Mass balance and boundedness.}
\beP[A positive mass balance gives a compact absorbing set][][compact]
For $\dot x=f(x)$ and $q\gg0$, suppose throughout the nonnegative orthant that
\[
 q^\top f(x)\leq a-bq^\top x,\qquad a\geq0,\quad b>0.
\]
Then $W(t)=q^\top x(t)$ satisfies
\[
 W(t)\leq a/b+(W(0)-a/b)e^{-bt}.
\]
All solutions are forward complete. Every set
$\{x\geq0:q^\top x\leq L\}$ with $L>a/b$ is compact,
positively invariant, and absorbs every trajectory in finite time.
\eeP
\begin{proof}
Multiply the differential inequality for $W-a/b$ by $e^{bt}$ and
integrate. At $W=L>a/b$ the derivative is strictly negative, proving
invariance, and the displayed bound gives absorption. Positivity of
all components of $q$ makes the sublevel set bounded. The estimate keeps
each solution in a compact set, so local existence extends to all forward times.
\end{proof}
This is the standard dissipative mass-balance estimate behind compact
population models and persistence arguments \cite{SmithThieme2011,HS10}.
\fi
\paragraph{Positive weights: one test for stability, decay and comparison.}

\beP[Collatz--Wielandt bounds for Metzler matrices][][ColWie]
For a Metzler matrix $A$, the following are equivalent:
\[
 s(A)<0;\qquad \exists v\gg0:\ Av\ll0;
 \qquad \exists q\gg0:\ q^\top A\ll0.
\]
For every $v\gg0$,
\[
 \min_i\frac{(Av)_i}{v_i}\leq s(A)
 \leq\max_i\frac{(Av)_i}{v_i}.
\]
If a nonnegative solution satisfies
$\dot z\leq Az$ componentwise, and $q^\top A\leq-\alpha q^\top$
for some $\alpha>0$, then
\[
 q^\top z(t)\leq e^{-\alpha t}q^\top z(0).
\]
\eeP
\begin{proof}
If $A$ is Hurwitz, its positive semigroup gives
$(-A)^{-1}=\int_0^\infty e^{At}\,dt\geq0$.
Choose $v=(-A)^{-1}\mathbf1$ and
$q^\top=\mathbf1^\top(-A)^{-1}$; their components are strictly
positive and the strict inequalities follow. Conversely, a positive
left or right weight with strict decrease makes the associated
weighted column or row sums strictly negative, and bounds $s(A)<0$.
The two-sided bounds follow by diagonally scaling with $v$ and
comparing the row sums of a shifted nonnegative matrix.
The final assertion follows from
$\frac{d}{dt}(q^\top z)\leq-\alpha q^\top z$ and Gronwall's inequality.
\end{proof}
These are classical nonsingular $M$-matrix and positive-systems
criteria \cite{Berman}, \cite{rantzer2015scalable}. A uniform comparison on a forward-invariant
region upgrades a local invasion test to an extinction estimate on
that region. Without the comparison hypothesis, a stable boundary
equilibrium alone says nothing about large introductions.

\paragraph{Comparison matrices and explicit extinction rates.}
The preceding positive-weight criterion becomes an explicit decay test
when the comparison matrix is written as $C-V$. The next theorem applies
to any nonnegative differential subsolution; a siphon supplies one
important application.

\beT[A Metzler comparison inequality gives explicit exponential decay]
[comparison decay][Cmp]
\emph{Scope: any nonnegative absolutely continuous function satisfying the
differential inequality; no condition is placed on the other variables.}

Let $V\in\mathbb R^{n\times n}$ be a nonsingular $M$-matrix and $C\ge0$ with $\rho(CV^{-1})<1$.
Then $V-C$ is a nonsingular $M$-matrix. Put
\[
 p=(V-C)^{-1}\mathbf1\gg0,\qquad c=\frac1{\max_ip_i},\qquad
 w^\top=\mathbf1^\top(V-C)^{-1}\gg0,\qquad c'=\frac1{\max_iw_i}.
\]
If $z:[0,\infty)\to\mathbb R_+^n$ is locally absolutely continuous and
$\dot z\le(C-V)z$ almost everywhere, componentwise, then
\[
 \max_i\frac{z_i(t)}{p_i}\le e^{-ct}\max_i\frac{z_i(0)}{p_i},\qquad
 L=w^\top z\ \text{ satisfies }\ \dot L\le-c'L\quad\text{a.e.}
\]
In particular $z(t)\to0$ exponentially.
\eeT
\begin{proof}
$(V-C)^{-1}=V^{-1}\sum_{k\ge0}(CV^{-1})^k\ge V^{-1}\ge0$ and $V-C$ is a
$Z$-matrix, so it is a nonsingular $M$-matrix; $V^{-1}$ has no zero row or
column, hence $p,w\gg0$. The matrix $C-V$ is Metzler and $(C-V)p=-\mathbf1\le-cp$.
The function $r(t)=\max_i z_i(t)/p_i$ is locally absolutely continuous.
At almost every time, its derivative is the maximum of
$\dot z_i(t)/p_i$ over the indices attaining the maximum. For each such
index, $z_i=rp_i$, $z_j\le rp_j$, and the Metzler property gives
\[
 \dot z_i\le((C-V)z)_i\le r((C-V)p)_i=-r\le-crp_i.
\]
Thus $\dot r\le-cr$ almost everywhere, and integration proves the
weighted maximum estimate. For the functional, $w\ge0$ gives
$\dot L\le w^\top(C-V)z=-\mathbf1^\top z\le-L/\max_iw_i$.
\end{proof}


For a siphon, apply the theorem to $z=x_\Sigma$ after verifying the
comparison inequality throughout the stated invariant region. No
cooperativity or convergence assumption on the resident variables is
needed. The inequality may also depend on time, provided the same
constant matrices $C,V$ bound it almost everywhere. A stable transverse
Jacobian on the face alone does not provide this global comparison.

\paragraph{What the siphon lattice does, and does not, tell the linear graph.}
At a relative-interior point of $\mathcal F_\Sigma$, let $G_\Sigma$ have the
directed edge $j\to i$ for $i\ne j$ when
$(\JtrS{\Sigma})_{ij}>0$.

\beT[Siphons constrain the linear influx graph][siphons and the influx graph][Graph]\label{p:toolGraph}\label{t:Sip}
For mass-action kinetics with positive rate constants, if
$T\subseteq\Sigma$ is a siphon, then $T$ is predecessor-closed in
$G_\Sigma$: $i\in T$ and $j\to i$ imply $j\in T$.
If every reaction producing a species of $\Sigma$ has exactly one
reactant molecule in $\Sigma$, the converse also holds.
Thus, in this first-order setting, minimal nonempty siphons contained
in $\Sigma$ are the source strongly connected components of
$G_\Sigma$. They need not be all its strongly connected components.
\eeT
\begin{proof}
An off-diagonal edge arises from a reaction of order one in $\Sigma$:
its unique reactant in $\Sigma$ is $j$, and it produces $i\ne j$.
If $i\in T$, the siphon property forces that reactant $j$ into $T$.
For the converse, take a reaction producing $i\in T$. Its unique
$\Sigma$ reactant is $j$. If $j=i$, it consumes $T$; otherwise it
gives the edge $j\to i$ and predecessor closure forces $j\in T$.
The last assertion is the source-component characterization of
minimal nonempty predecessor-closed sets.
\end{proof}
This is the elementary graph reduction behind monomolecular siphon
analysis \cite{AdLS,ShiuStu,Berman}. \label{pg:2A2B}The added first-order hypothesis is
essential. The networks
\[
 A\to B,\quad B\to A,
 \qquad\text{and}\qquad
 2A\to B,\quad 2B\to A
\]
have the same sole nonempty siphon $\{A,B\}$, but at its zero face the
first transverse graph is strongly connected and the second has no
edges. Adding $A\to\varnothing$, $B\to\varnothing$ to both preserves
the example and gives negative diagonal entries in the second.
Hence the siphon lattice alone cannot determine the linear influx
graph without information about reaction order. Moreover, when the
transverse graph is reducible,
$s(\JtrS{\Sigma})=\max_B s(J_B)$ in Frobenius form, and a
supercritical block may fail to invade from introductions confined
to components which cannot reach it.

% Exit test: unused by the short paper, kept for the long versions.
\ifshort\else
\paragraph{Biologically labelled splittings and the exit test.}
If a biological convention has already specified a nonnegative matrix
$F$ of new infections, put $V=F-\JtrS{\Sigma}$. This is a meaningful
splitting only after its matrix conditions have been checked. The
labels selecting $F$ are extra model data; stoichiometry alone need
not distinguish a newly infected individual from an internal transfer.

\beT[When the labelled loss matrix is admissible]\label{p:toolExit}
For $F\geq0$, the matrix $V=F-\JtrS{\Sigma}$ is a nonsingular
$M$-matrix if and only if
\[
 V_{ij}\leq0\quad(i\ne j),\qquad
 \exists q\gg0:\ q^\top V\gg0.
\]
In the special transfer-and-exit representation
\[
 Q=-V,\quad Q_{ij}=t_{ij}\geq0\ (i\ne j),\quad
 Q_{jj}=-d_j-\sum_{i\ne j}t_{ij},\qquad d_j\geq0,
\]
this holds exactly when every compartment can reach one with a
strictly positive exit rate $d_j$ along transfer edges $j\to i$.
Equivalently, every terminal strongly connected component contains
a positive exit.
\eeT
\begin{proof}
The first equivalence is the preceding Metzler criterion applied to
$-V$. For the second, $Q$ is the generator, in column convention, of
a finite transfer process killed at rates $d_j$. If every state can
reach an exit, choose a common finite time and a positive lower bound
on the exit probability from every initial state. Iterating gives
exponential decay of the surviving mass, hence $Q$ is Hurwitz.
Conversely a terminal component without exit conserves its mass;
its generator has eigenvalue zero, so $Q$ cannot be Hurwitz.
\end{proof}
This is a standard compartmental, or $M$-matrix, test
\cite{Berman,Van}.
It turns the question of an epidemiologically meaningful splitting
into a labelled modelling choice followed by a precise admissibility
test.
\fi

\paragraph{From local invasion to persistence.}
The appropriate resident object is an equilibrium, periodic orbit, or
invariant measure, not necessarily a fixed point. Conservation can
exclude some boundary faces outright; positive invasion weights can
repel others. Uniform persistence additionally needs a compact
absorbing setting and control of all relevant boundary recurrent
sets, as provided by an acyclic Morse decomposition or a suitable
average-Lyapunov criterion \cite{FRT94,HS10,HofSch,SmithThieme2011}. A list of
positive invasion numbers at selected equilibria is not a substitute
for checking those hypotheses. 